\section{Conclusion}
\label{sec:conclusion}

We asked whether an agent's written prediction controls its next action when the knowledge it states has been learned into the weights, and what in the training data decides this. In \alfcf{}, where object locations can only be learned from fine-tuning, the answer depends on the data. Qwen3-4B agents fine-tuned on data in which every prediction is followed copy their prediction into the first search step, although deleting the prediction changes almost nothing. When a quarter or half of the training predictions are contradicted by the action, the agents keep writing predictions but follow a replaced one in at most about 1\% of first steps, as an imitator of such data would.

Two implications follow for work that reads or edits agents' written reasoning. First, deletion and editing answer different questions. Because our data also teach acting without a prediction, a deletion test would have called the predictions of the always-followed agents decorative, while writing another object's prediction at every step steers them through the whole episode. Second, in our setting the training data decided whether a written prediction controlled the decision. Agents whose training predictions were always borne out by the next action followed a replaced prediction at about 80\% of first steps and a randomly ordered one at 90\%. Contradicting a quarter of the training predictions removed this, as the data imply, while the agents kept writing predictions.
